\usepackage{xcolor}
\usepackage{afterpage}
\usepackage{pifont,mdframed}
\usepackage[bottom]{footmisc}

\newcommand{\inputfile}{\texttt{stdin}}
\newcommand{\outputfile}{\texttt{stdout}}
\makeatletter
\renewcommand{\this@inputfilename}{\texttt{stdin}}
\renewcommand{\this@outputfilename}{\texttt{stdout}}
\makeatother

È l'ora di pranzo allo stage per le Olimpiadi di Informatica e tutti gli studenti sono affamati! Tommaso aveva promesso di preparare la pizza per tutti, ma come al solito è troppo lento. Samuele, stufo di aspettare, decide di andare lui stesso in pizzeria a prendere le pizze.

Via Gemona, la via dove si trova il Toppo, è composta da $N$ luoghi ($2 \le N \le 100\,000$), numerati da $1$ a $N$, disposti in fila lungo la strada. I luoghi sono collegati in sequenza: il luogo $1$ è collegato al luogo $2$, il luogo $2$ è collegato al luogo $3$, e così via fino al luogo $N$.

Ci sono $N-1$ partecipanti allo stage, uno per ogni luogo da $1$ a $N-1$. Il Toppo (la sede dello stage dove tutti devono tornare) si trova al luogo $N$. Ogni partecipante deve tornare al Toppo partendo dal proprio luogo.

Tra un luogo e il successivo c'è una certa distanza (in minuti di cammino). La distanza tra il luogo $i$ e il luogo $i+1$ è indicata con $t_i$ minuti.

Lungo Via Gemona, in $K$ dei luoghi ($1 \le K \le N$) ci sono dei chioschi che vendono tranci di pizza. Il chiosco nel luogo $i$ vende un trancio di pizza con un livello di soddisfazione $y_i$ (misurato in "quanto vale la pena fermarsi").

Ogni partecipante è disposto a fermarsi al massimo in un chiosco lungo il suo percorso verso il Toppo, ma solo se il tempo extra aggiunto al suo percorso è al massimo uguale al livello di soddisfazione della pizza che prenderebbe.

Nota importante: Un partecipante può decidere di tornare indietro per prendere un trancio di pizza e poi riprendere il cammino verso il Toppo, se questo ha senso dal punto di vista del tempo extra impiegato.

Esempio: Se un partecipante si trova al luogo $5$ e il Toppo è al luogo $7$, il percorso più breve è $5 \to 6 \to 7$. Tuttavia, se c'è un chiosco al luogo $3$ con una pizza molto soddisfacente, il partecipante potrebbe decidere di fare: $5 \to 4 \to 3$ (prende la pizza) $\to 4 \to 5 \to 6 \to 7$.

\InputFile
La prima riga contiene due numeri interi $N$ e $K$ separati da spazi.

Ognuna delle successive $N-1$ righe contiene un intero $t_i$ che rappresenta i minuti necessari per andare dal luogo $i$ al luogo $i+1$ ($1 \le t_i \le 10^4$).

Le successive $K$ righe descrivono ciascuna un chiosco di pizza con due interi: l'indice del suo luogo e il suo livello di soddisfazione (un intero positivo al massimo $10^9$). Più chioschi possono trovarsi nello stesso luogo.

\OutputFile
L'output deve consistere di $N-1$ righe. La riga $i$ contiene il singolo intero $1$ se il partecipante nel luogo $i$ può fermarsi a mangiare una pizza lungo il percorso verso il Toppo, e $0$ altrimenti.

\Constraints
\begin{itemize}
    \item $2 \le N \le 100\,000$.
    \item $1 \le K \le N$.
    \item $1 \le t_i \le 10^4$.
    \item $1 \le y_i \le 10^9$.
\end{itemize}

\Scoring
\begin{itemize}
    \item \subtask Casi d'esempio.
    \item \subtask $N \le 10$, $K \le 10$.
    \item \subtask $N \le 1\,000$, $K \le 1\,000$.
    \item \subtask $K = 1$.
    \item \subtask $t_i = 1$ per ogni strada.
    \item \subtask Nessuna limitazione aggiuntiva.
\end{itemize}

\Examples
\begin{example}
    \exmpfile{input0.txt}{output0.txt}
\end{example}

\Explanations
La disposizione dei luoghi lungo Via Gemona è: $1 \xrightarrow{10} 2 \xrightarrow{5} 3 \xrightarrow{8} 4 \xrightarrow{3} 5$ (Toppo).
Ci sono chioschi:
\begin{itemize}
    \item Nel luogo $2$ con soddisfazione $15$.
    \item Nel luogo $4$ con soddisfazione $6$.
\end{itemize}

Partecipante 1:
\begin{itemize}
    \item Percorso diretto: $1 \to 2 \to 3 \to 4 \to 5$ (tempo: $10+5+8+3 = 26$ minuti).
    \item Con pizza al luogo $2$: $1 \to 2$ (prende la pizza) $\to 3 \to 4 \to 5$ (tempo: $26$ minuti).
    \item Tempo extra: $0$, soddisfazione: $15 \implies$ può fermarsi! ($0 \le 15$) \checkmark
\end{itemize}

Partecipante 2:
\begin{itemize}
    \item Percorso diretto: $2 \to 3 \to 4 \to 5$ (tempo: $5+8+3 = 16$ minuti).
    \item Con pizza al luogo $2$: già lì! (tempo: $16$ minuti).
    \item Tempo extra: $0 \implies$ può fermarsi! \checkmark
\end{itemize}

Partecipante 3:
\begin{itemize}
    \item Percorso diretto: $3 \to 4 \to 5$ (tempo: $8+3 = 11$ minuti).
    \item Con pizza al luogo $4$: $3 \to 4$ (prende la pizza) $\to 5$ (tempo: $11$ minuti).
    \item Tempo extra: $0$, soddisfazione: $6 \implies$ può fermarsi! ($0 \le 6$) \checkmark
\end{itemize}

Partecipante 4:
\begin{itemize}
    \item Percorso diretto: $4 \to 5$ (tempo: $3$ minuti).
    \item Con pizza al luogo $4$: già lì! (tempo: $3$ minuti).
    \item Tempo extra: $0 \implies$ può fermarsi! \checkmark
\end{itemize}
